% ==============================================================================
% Esqueleto do artigo — Zero-Shot Cross-Lingual Audio Deepfake Detection
% (BRSpeech-DF, pt-BR) — alvo: DFF-2026 (workshop @ ACM MM 2026, Rio de Janeiro)
%
% NOTA: o template oficial do DFF-2026 ainda não foi publicado (página fora do
% ar no momento em que isso foi escrito). Este preâmbulo usa o formato padrão
% "sigconf" do acmart, que é o mais comum em workshops ACM. Quando o CfP oficial
% sair, troque \acmConference{}{}{} pelos dados reais e confira se pedem alguma
% opção extra na classe (ex.: review anônima -> adicionar "anonymous" e "review").
% ==============================================================================
\documentclass[sigconf]{acmart}

% Evita a caixa de referência bibliográfica completa no rodapé da 1a página
% (padrão em vários workshops ACM; remova se o CfP pedir o formato completo)
\settopmatter{printacmref=false}
\renewcommand\footnotetextcopyrightpermission[1]{}

\usepackage{booktabs}   % tabelas (EER/MCC/Acc, robustez)
\usepackage{graphicx}   % figuras (curvas de robustez, GATR, oclusão)

\begin{document}

% ------------------------------------------------------------------------
% TODO: título de trabalho -- ajustar depois que os resultados finais saírem
% ------------------------------------------------------------------------
\title{Zero-Shot Cross-Lingual Audio Deepfake Detection on Brazilian Portuguese Speech}

% TODO: nome completo, afiliação e e-mail de cada autor
\author{João Victor}
\affiliation{%
  \institution{Universidade Federal de Goiás (UFG) / CEIA}
  \city{Goiânia}
  \country{Brazil}
}
\email{TODO@ufg.br}

\renewcommand{\shortauthors}{João Victor et al.}

\begin{abstract}
% TODO (~150-200 palavras): motivação (escassez de detectores nativos em
% pt-BR), o que foi feito (avaliação zero-shot cross-lingual de um detector
% wav2vec 2.0 / XLS-R pré-treinado, sem fine-tuning, no BRSpeech-DF), como foi
% avaliado (EER, MCC, Acurácia; robustez sob ruído/MP3/G.711; explicabilidade
% no domínio do tempo via GATR e oclusão), e o achado principal (preencher
% depois de rodar a pipeline no ambiente final).
TODO: abstract.
\end{abstract}

\keywords{audio deepfake detection, anti-spoofing, cross-lingual, wav2vec 2.0, XLS-R, explainability, Portuguese}

\maketitle

% ==============================================================================
\section{Introduction}
% Notebook: célula 0 (markdown de abertura) tem a motivação em pt-BR --
% traduzir e expandir, não copiar literal.
% Pontos a cobrir:
%  - por que deepfake de voz em pt-BR é pouco coberto na literatura
%  - por que testar zero-shot antes de investir em fine-tuning é uma pergunta
%    de pesquisa válida (custo, disponibilidade de dados rotulados em pt-BR)
%  - contribuições do artigo (bullet list de 3-4 itens)
TODO.

% ==============================================================================
\section{Related Work}
% Os 4 papers-base (ver references.bib) + posicionar frente à revisão
% sistemática PRISMA (ENIAC 2026) já feita por você -- deixar claro que este
% artigo é um estudo empírico complementar, não uma sobreposição.
\subsection{Self-Supervised Front-Ends for Anti-Spoofing}
\citet{tak2022automatic} TODO.

\subsection{Robustness and Real-World Generalization}
\citet{channing2024toward} TODO.

\subsection{Explainability in Audio Deepfake Detection}
\citet{grinberg2025does} and \citet{fullwood2025seeing} TODO.

% ==============================================================================
\section{Methodology}
% Notebook, seções 2-3 (pré-processamento + modelo). Descrever como pipeline,
% não como "célula X faz Y" -- traduzir para prosa científica.
\subsection{Pre-processing}
% mono -> resample 16 kHz -> janela fixa de 64.600 amostras (~4.04s) ->
% normalização (Tak et al., 2022)
TODO.

\subsection{Model}
% front-end SSL wav2vec 2.0 / XLS-R + back-end (pooling global + linear),
% checkpoint pré-treinado (nii-yamagishilab/mms-300m-anti-deepfake), sem
% fine-tuning -- deixar claro que é zero-shot cross-lingual (treinado
% majoritariamente em inglês/ASVspoof, aplicado direto em pt-BR)
TODO.

% ==============================================================================
\section{Experimental Setup}
% Notebook, seção 4 (dados) + CFG. Vira uma tabela de configuração.
\begin{table}[t]
  \caption{Experimental configuration.}
  \label{tab:config}
  \begin{tabular}{ll}
    \toprule
    Parameter & Value \\
    \midrule
    Sample rate & 16{,}000 Hz \\
    Fixed window & 64{,}600 samples ($\approx$4.04 s) \\
    Checkpoint & \texttt{nii-yamagishilab/mms-300m-anti-deepfake} \\
    Dataset & AKCIT-Deepfake/BRSpeech-DF (test split) \\
    Samples per class & 200 (balanced) \\
    Seed & 42 \\
    \bottomrule
  \end{tabular}
\end{table}
TODO: descrever o BRSpeech-DF (tamanho, origem, como o rótulo bonafide/spoof
é resolvido), hardware usado, e re-confirmar que a versão final rodou no
ambiente único definido (ver pendência nas instruções do projeto).

% ==============================================================================
\section{Results}
% Notebook, seções 5-7 (inferência, métricas, robustez)
\subsection{Zero-Shot Performance}
% Tabela EER / MCC / Acurácia no áudio íntegro + matriz de confusão
TODO.

\subsection{Robustness Under Real-World Degradations}
% Tabela/gráfico comparando áudio íntegro vs. ruído gaussiano vs. MP3 vs. G.711
TODO.

% ==============================================================================
\section{Explainability}
% Notebook, seções 8-9 (GATR + oclusão)
\subsection{GATR: Time-Domain Relevancy}
TODO.

\subsection{Occlusion Analysis: Speech vs.\ Silence}
% inclui o achado da agregação via VAD (fala vs. silêncio)
TODO.

% ==============================================================================
\section{Discussion and Limitations}
% janela fixa de 4s, tamanho de amostra (n_per_class=200 -- variância do EER),
% generalização (um único checkpoint, um único dataset pt-BR), decisão sobre
% incluir ou não W2V2-AASIST como baseline
TODO.

% ==============================================================================
\section{Conclusion and Future Work}
TODO.

% ==============================================================================
\bibliographystyle{ACM-Reference-Format}
\bibliography{references}

\end{document}
